%% kapa: exam=1 round=28 year=2017 period=2 form=A subject=accounting subject_ko=회계학 seq=35 number=75
%% source: Q-Net 감정평가사 기출문제 (cst00309) artlSeq=5199603 file=extracted/1008/5199603/2017년 제28회 감정평가사 1차시험 문제지/2교시 A형.pdf
%% conversion: kordoc 4.24.51 -> tools/gen1.py (auto); review=proofread-sonnet; flags=none
\begin{question}{75}
(주)대한은 X, Y, Z 제품을 생산ㆍ판매하고 있으며, 20x1년도 제품별 예산손익계산서는 다음과 같다.
\begin{framedtext}
\noindent\begin{tabular}{@{}l@{\hspace{1.2em}}l@{\hspace{2em}}r@{\hspace{2em}}r@{\hspace{2em}}r@{}}
\multicolumn{2}{c}{구분} & \multicolumn{1}{c}{X제품} & \multicolumn{1}{c}{Y제품} & \multicolumn{1}{c}{Z제품} \\ \hline
매출액 & & ￦100,000 & ￦200,000 & ￦150,000 \\
매출원가: & 변동원가 & 40,000 & 80,000 & 60,000 \\
 & 고정원가 & 30,000 & 70,000 & 50,000 \\ \cline{3-5}
매출총이익 & & ￦30,000 & ￦50,000 & ￦40,000 \\
판매관리비: & 변동원가 & 20,000 & 10,000 & 10,000 \\
 & 고정원가 & 20,000 & 20,000 & 20,000 \\ \cline{3-5}
영업이익(손실) & & (￦10,000) & ￦20,000 & ￦10,000 \\
\end{tabular}
\end{framedtext}
(주)대한의 경영자는 영업손실을 초래하고 있는 X제품의 생산을 중단하려고 한다. X제품의 생산을 중단하면, X제품의 변동원가를 절감하고, 매출원가에 포함된 고정원가의 40\%와 판매관리비에 포함된 고정원가의 60\%를 회피할 수 있다. 또한, 생산중단에 따른 여유생산능력을 임대하여 ￦10,000의 임대수익을 얻을 수 있다. X제품의 생산을 중단할 경우, 20x1년도 회사 전체의 예산 영업이익은 얼마나 증가(또는 감소)하는가? (단, 기초 및 기말 재고자산은 없다.) \par
\begin{choices}
  \item ￦4,000 감소
  \item ￦5,000 증가
  \item ￦6,000 감소
  \item ￦7,000 증가
  \item ￦8,000 증가
\end{choices}
\end{question}
\answer{3}
